\documentclass{scrartcl}
\usepackage{graphicx}  % required for inserting images

% for cyrillic and Bulgarian language
\usepackage[utf8]{inputenc}
\usepackage[T2A]{fontenc}
\usepackage[bulgarian]{babel}

% math packages
\usepackage{amsmath}
\usepackage{float}
\usepackage{amsfonts}
\usepackage{amssymb}
\usepackage{amsthm}
\usepackage{cancel}

\usepackage{ragged2e}  % adds FlushLeft

\usepackage{listings}
\usepackage{xcolor}

\definecolor{codegreen}{rgb}{0,0.6,0}
\definecolor{codegray}{rgb}{0.5,0.5,0.5}
\definecolor{codepurple}{rgb}{0.58,0,0.82}
\definecolor{backcolour}{rgb}{0.95,0.95,0.92}

\lstdefinestyle{mystyle}{
    % caption=Example C++,
    label={lst:listing-cpp},
    language=C++,
    backgroundcolor=\color{backcolour},   
    commentstyle=\color{codegreen},
    keywordstyle=\color{magenta},
    numberstyle=\tiny\color{codegray},
    stringstyle=\color{codepurple},
    basicstyle=\ttfamily\footnotesize,
    breakatwhitespace=false,         
    breaklines=true,                 
    keepspaces=true,                 
    numbers=left,       
    numbersep=5pt,                  
    showspaces=false,                
    showstringspaces=false,
    showtabs=false,                  
    tabsize=2,
}
\lstset{style=mystyle}

\usepackage[colorlinks=true, linkcolor=blue, urlcolor=cyan]{hyperref}

\newtheorem{theorem}{Теорема}
\newtheorem{definition}{Дефиниция}

\title{Увод в програмирането - основни конструкции}
\author{Ивайло Андреев + gemini 3 Flash, chat GPT 5.4}
\date{16 май 2026}

\begin{document}
\maketitle

\tableofcontents

\newpage

\begin{FlushLeft}

\section{Принципи на структурното програмиране}
Процедурната парадигма служи като основа за разбирането на съвременните подходи и е неизменна част от изучаването на алгоритми и структури от данни. Програмирането може да бъде разделено на две основни парадигми:
\begin{itemize}
    \item \textbf{Императивно програмиране:} Описва \textit{как} се изпълнява дадена задача, чрез поредица от инструкции, които променят състоянието на програмата (например чрез променливи и команди за контрол на потока). То включва процедурното и обектно-ориентираното програмиране.
    \item \textbf{Декларативно програмиране:} Описва \textit{какво} трябва да се постигне, без изрично да се указват стъпките за изпълнение (например функционално програмиране и логическо програмиране).
\end{itemize}

\textbf{Фундаментална теорема на структурното програмиране (Теорема на Бьом-Якопини):} Всеки алгоритъм може да бъде изразен чрез използването на само три основни управляващи конструкции:
\begin{enumerate}
    \item \textbf{Последователност (Sequence):} Инструкциите се изпълняват една след друга в реда, в който са написани.
    \item \textbf{Разклонение (Selection/Condition):} Избор между два или повече пътя на изпълнение в зависимост от логическо условие.
    \item \textbf{Цикъл (Iteration/Loop):} Многократно изпълнение на поредица от инструкции, докато дадено условие е изпълнено.
\end{enumerate}
Основната цел на структурното програмиране е да направи програмите по-разбираеми, по-лесни за проследяване и модификация, като се избягва използването на неструктурирани преходи като \texttt{goto}.

\section{Управление на изчислителния процес}
Управлението на изчислителния процес се осъществява чрез инструкции, които определят реда на изпълнение на действията. Програмата се състои от изрази (които пресмятат и връщат стойност) и оператори (които извършват действия). Блок от код се огражда с къдрави скоби \texttt{\{\}}.

\subsection{Условни оператори}
Условните оператори позволяват разклоняване на логиката. Използват се \texttt{if}, \texttt{if-else} и \texttt{switch} за многозначно разклоняване.
\begin{lstlisting}
// Example of conditional statements and basic syntax
int x = 10;
if (x > 5) {
    // Executed if the condition is true
    x--;
} else {
    // Executed if the condition is false
    x++;
}

// Inline ternary conditional choice
int y = (x == 9) ? 1 : 0; 

// Switch selector statement
switch (y) {
    case 0:
        // Handle case 0 block
        break; // Stop fall-through
    case 1:
        // Handle case 1 block
        break;
    default:
        // Optional catch-all block
        break;
}
\end{lstlisting}

\subsection{Оператори за цикъл}
Циклите осигуряват многократно изпълнение на код. Те се делят на цикли с предусловие (\texttt{while}), цикли с постусловие (\texttt{do-while}) и цикли с брояч (\texttt{for}).
\begin{lstlisting}
// 1. Pre-condition loop (evaluates condition before running body)
int i = 0;
while (i < 5) {
    // Loop code statements
    i++;
}

// 2. Post-condition loop (guaranteed to execute at least once)
int j = 0;
do {
    // Executes body, then evaluates condition
    j++;
} while (j < 5);

// 3. Counter-controlled loop iteration
for (int k = 0; k < 5; k++) {
    // Initialization, condition check, and step execution
}
\end{lstlisting}

\section{Променливи}
Променливата е наименувана област от паметта, която съхранява стойност от определен тип. Всяка променлива има адрес, съответстващ на позицията ѝ в паметта.

\subsection{Локални и глобални променливи}
Основните разлики между променливите се свеждат до техния обхват (видимост) и време на живот:
\begin{itemize}
    \item \textbf{Локални променливи:} Дефинират се вътре във функция или блок \texttt{\{\}}. Те се създават при достигане на дефиницията им и се унищожават при излизане от блока (разполагат се в \textbf{Стека}). Видими са само в този блок.
    \item \textbf{Глобални променливи:} Дефинират се извън всички функции, в началото на файла. Те съществуват през цялото време на изпълнение на програмата (разположени в \textbf{Data/BSS сегмента}) и са достъпни от всяка функция в рамките на файла.
\end{itemize}

\subsection{Инициализация и оператор за присвояване}
\begin{itemize}
    \item \textbf{Оператор за присвояване (\texttt{=}):} Променя текущата стойност на вече съществуваща променлива, като оценява израза отдясно и го записва в паметта на променливата отляво.
    \item \textbf{Инициализация:} Задаване на начална стойност на променливата в момента на нейното деклариране. Важно е да се отбележи, че неинициализираните локални променливи съдържат произволни данни от паметта (т.нар. "боклук"), което води до грешки и недефинирано поведение, докато глобалните се инициализират автоматично с 0.
\end{itemize}

\section{Функции и процедури}
Функцията е самостоятелен блок от код, който решава подзадача и може да бъде извикан многократно. В C++ процедурите се представят като функции с тип на връщане \texttt{void}.
\begin{itemize}
    \item \textbf{Формални параметри:} Променливите, изброени в заглавието на функцията.
    \item \textbf{Фактически параметри (Аргументи):} Реалните стойности или променливи, които се подават в скобите при извикването.
\end{itemize}

\subsection{Предаване на параметри}
По официален конспект се разглеждат следните основни механизми за предаване:
\begin{itemize}
    \item \textbf{Предаване по стойност (Pass by Value):} Създава се локално копие на аргумента в стековата рамка на функцията. Промените вътре във функцията не влияят на оригиналната променлива.
    \item \textbf{Предаване по име (Pass by Name):} Исторически модел (използван например в ALGOL 60), при който аргументът се подставя на местата на формалния параметър в тялото на функцията при извикване. Този механизъм не е част от стандартния C++, но се разглежда теоретично в конспекта.
\end{itemize}
\textit{Забележка:} За да се променят оригиналните стойности на променливите, в практиката на C++ се използва предаване чрез псевдоним (референция) или указател, но теоретичният акцент в академичния конспект пада върху горните два модела.
\begin{lstlisting}
// Demonstrating parameter passing models
void modifyValue(int num) {
    // Pass by value: changes only the local stack copy
    num = 100; 
}

void modifyByReference(int& num) {
    // Pass by reference: changes the original argument
    num = 100;
}

int main() {
    int x = 5;

    modifyValue(x);       // x remains 5
    modifyByReference(x); // x becomes 100

    return 0;
}
\end{lstlisting}

\subsection{Типове и проверка за съответствие на тип}
C++ е статично типизиран език. Компилаторът извършва строга проверка за съответствие на типовете (\textbf{Type Checking}) по време на компилация. Когато подадем аргумент на функция, компилаторът проверява дали типът му съвпада с дефинирания формален параметър. При несъвпадение се прави опит за автоматично (неявно) преобразуване на типа, а ако такова е невъзможно, компилацията се прекъсва с грешка за несъвместимост.

\end{FlushLeft}

\end{document}